% The conversion algorithm as a block diagram. One box per step, the key
% methods named inside it, and on the right the invariant each step is what
% protects -- the three requirements the transfer ledger checks afterwards.
\begin{figure}[!ht]
\centering
\begin{tikzpicture}[
  font=\footnotesize, >=Latex, node distance=3mm,
  blk/.style={draw, rounded corners=3pt, align=left, text width=63mm,
               inner sep=2mm},
  s1/.style={blk, fill=blue!8,   draw=blue!55},
  s2/.style={blk, fill=orange!9, draw=orange!65},
  s3/.style={blk, fill=teal!9,   draw=teal!60},
  s4/.style={blk, fill=violet!9, draw=violet!55},
  inv/.style={draw, rounded corners=2pt, align=left, text width=25mm,
              inner sep=1.5mm, fill=black!4, draw=black!35,
              font=\scriptsize},
  hdr/.style={font=\footnotesize\bfseries},
  ar/.style={->, thick, black!60},
  tie/.style={->, black!35, dashed, shorten >=1pt},
]

\node[s1] (v) {%
  \textbf{1\quad Nodal volumes}\\[2pt]
  restricted Voronoi on the CD node cloud\\
  Monte Carlo in the bounding box, half-space test $\bm N\xv\le\bm b$\\
  nearest node by $k$-d tree, $V_j=(M_j/M)V_\Omega$};

\node[s2, below=of v] (n) {%
  \textbf{2\quad How many loops}\\[2pt]
  $\lambda^k_{sj}=n^k_{sL}(\xv_j)V_j/\Omega$, summed to $\Lambda^k_s$\\
  $N_{\mathrm{draw}}=\operatorname{round}(\Lambda^k_s)$, not a Poisson total\\
  nodes drawn multinomially, once per family};

\node[s3, below=of n] (g) {%
  \textbf{3\quad Size and shape}\\[2pt]
  $m_\ell$ drawn from the home node's log-normal\\
  rescaled so $\sum_\ell m_\ell$ matches the content exactly\\
  equal-\emph{area} polygon on the habit plane, elliptical\\
  center by rejection sampling inside the crystal};

\node[s4, below=of g] (f) {%
  \textbf{4\quad What the domain refuses}\\[2pt]
  predict the solver's refusal of loops leaving the grain\\
  keep the refused share in the continuum\\
  count the refusal in defects, never in loops};

\draw[ar] (v) -- (n);
\draw[ar] (n) -- (g);
\draw[ar] (g) -- (f);

\node[inv, right=4mm of v] (i0) {%
  $\sum_j V_j=V_\Omega$\\ identically, however coarse the sampling};
\node[inv, right=4mm of n] (i1) {%
  \textbf{(I)} loop number,\\ exact to the integer draw};
\node[inv, right=4mm of g] (i2) {%
  \textbf{(II)} stored defects,\\ exact to $10^{-16}$};
\node[inv, right=4mm of f] (i3) {%
  \textbf{(III)} sink strength,\\ continuous to ${\sim}2\%$};

\draw[tie] (v) -- (i0);
\draw[tie] (n) -- (i1);
\draw[tie] (g) -- (i2);
\draw[tie] (f) -- (i3);

\node[anchor=north west, font=\scriptsize, text width=95mm, black!70]
  at ($(f.south west)+(0,-3mm)$)
  {The right-hand column is what each step exists to protect: the three
   requirements of \cref{sec:transfer}, checked on the transfer ledger
   afterwards rather than assumed here.};
\end{tikzpicture}
\caption{The continuum-to-discrete conversion, step by step, with the methods
each step uses and the invariant it protects. The order is not arbitrary: the
volumes must precede the counts, because sampling on the density alone would
over-populate a refined region in exact proportion to how well it was resolved;
and the count must precede the sizes, because it is the rounded family total,
not a per-node rounding, that recovers the continuum density in the mean.}
\label{fig:conversion}
\end{figure}
